\section{Transformer baseline architecture}
\label{sec:transformer-baseline}
To establish a robust performance baseline, we first implemented a standard fine-tuning approach, which represents the conventional state-of-the-art for supervised text classification. Given the multilingual nature of the dataset, we selected XLM-RoBERTa \cite{conneau_unsupervised_2020}, a powerful multilingual encoder, as the foundation for our model.

The core architecture consists of the pre-trained XLM-RoBERTa model followed by a linear classification head. For each input text, the model produces a contextualized embedding for the special [CLS] token \cite{devlin_bert_2018}. This embedding, which represents the aggregated meaning of the entire input sequence, is then passed to a single feed-forward layer that projects it into a logit vector. The dimension of this vector is equal to the total number of sub-narrative labels in the taxonomy. A sigmoid activation function is applied to these logits during inference to produce a probability between 0 and 1 for each label, allowing for multi-label predictions.

We experimented with three variations of this baseline, each designed to systematically address the specific challenges identified in our data analysis.

\subsection{Vanilla fine-tuning}

Our first implementation establishes a foundational baseline using a straightforward fine-tuning methodology. We use the XLM-RoBERTa (xlm-roberta-base) model for several reasons: the task is multilingual, and XLM-RoBERTa is pre-trained on a large multilingual corpus covering many languages, enabling robust, language-agnostic representations; moreover, as a RoBERTa-family model it benefits from improved pre-training (e.g., dynamic masking) compared to the original BERT, resulting in stronger contextual embeddings.

The architectural setup is standard: the pre-trained XLM-RoBERTa encoder is followed by a linear classification head. During a forward pass the final hidden state of the special [CLS] token is taken as an aggregate representation of the input and passed through the classification head to produce a logit vector whose dimensionality matches the number of sub-narrative labels. At inference time a sigmoid is applied to each logit to produce independent probabilities for multi-label prediction.

For this vanilla baseline we frame the problem as a set of independent binary classification tasks (one per sub-narrative) and optimize using Binary Cross-Entropy with Logits Loss (BCEWithLogitsLoss). While simple and computationally efficient, this formulation assumes label independence and therefore does not capture structured co-occurrence patterns or hierarchical parent–child relationships present in the taxonomy; as our later variations show, explicitly addressing these properties can materially improve performance on tail and dependent labels.

\medskip
\noindent\textbf{Anticipated limitations.} Based on our data analysis, we anticipated that this vanilla approach would face significant limitations. The primary concern is its susceptibility to the severe class imbalance. By optimizing for overall accuracy, the standard BCE loss incentivizes the model to achieve high performance on the high-frequency ``head'' classes, often by simply predicting ``negative'' for the vast majority of rare ``tail'' classes.
This behavior is expected to result in poor performance on metrics sensitive to class imbalance, such as the macro-averaged F1 score, and serves as the primary motivation for the more advanced baselines and zero-shot architectures explored subsequently.

\subsection{Mitigating the class imbalance}
To directly address the class imbalance issue, our second variation enhances the loss function with positive class weighting. This technique modifies the standard BCE loss by increasing the penalty for misclassifying a positive instance of a rare class. For each class $c$ we compute a per-class weight \(\mathrm{pos\_weight}_c\) as the ratio of negative to positive training counts. Let $N_c$ and $P_c$ denote the number of negative and positive training examples for class $c$, respectively:
\[
	\mathrm{pos\_weight}_c \;=\; \frac{N_c}{P_c}.
\]

This weight is then integrated directly into the BCE loss calculation, effectively increasing the cost of a false negative for rare classes. For a common label where positive and negative samples are balanced, the weight is close to 1, behaving like the standard BCE loss. However, for a rare label with a \textbf{pos\_weight} of, for example, 65, the loss incurred for failing to predict that label when it is present is magnified 65 times.


This ensures that the model is more heavily penalized for failing to identify a rare narrative than for misclassifying a common one. In practice we pass the vector of per-class positive weights to the BCEWithLogitsLoss (or equivalent) implementation so rare classes contribute proportionally more to the loss. The approach is intended to improve recall on tail classes and, consequently, boost Macro F1 by forcing the model to pay closer attention to underrepresented labels.
